\scanpage{129}
\mspage{128}
\mspage{467}
$({}^{106}_{1})$ (3 octobre) Ni moi, ni Deligne n'avons jamais eu le moindre doute que les
conjectures de Weil puissent ne pas être valables, et je ne me rappelle pas \unsure{avoi|avoir}
entendu quiconque exprimer de tels doutes. Qualifier le ``résultat'' (i.e. la dé-\newline
monstration de ces conjectures) comme ``surprenant'', témoigne encore du propos
délibéré d'épater la galerie. D'ailleurs à aucun moment depuis l'introduction
de la ``topologie'' et de la cohomologie étales, je n'ai eu le sentiment que ces
conjectures étaient hors d'atteinte, mais plutôt \inserted{(}à partir de 1963\inserted{)} qu'elles ne
manqueraient pas d'être démontrées dans les toutes prochaines années. Au moment
de mon départ, en 1970, je n'avais guère de doute que Deligne, qui était le \unsure{mie|mieux}
placé de tous pour cela, ne tarderait pas à les prouver (ce qu'il n'a pas manqué
de faire), en même temps que les``conjectures standard sur les cycles algébriques
plus fortes, (qu'il s'est par contre attaché à discréditer).

C'est d'ailleurs avec raison que Deligne émet des réserves sur la validité
de ces dernières conjectures, dont je ne suis pas plus convaincu que lui. Mais
la portée d'une conjecture ne dépend pas du fait si elle finira par se révéler
vrai\inserted{e}, ou fausse, pas plus que son caractère de soi-disante ``difficulté'', qui la
rendrait``hors de portée'' - caractère entièrement subjectif, lui. Elle dépend \unsure{un|uni-}\newline
quement du fait si la\quad q u e s t i o n\quad sur laquelle la conjecture met le \unsure{doig|doigt}
(et qui n'avait pas été perçue, avant qu'elle ne soit posée) - si cette \unsure{questio|question}
touche à quelque chose de vraiment essentiel pour notre connaissance des choses.
Or il est evident (pour moi tout au moins!) qu'il ne saurait être question d'\unsure{avo|avoir}
une bonne compréhension des cycles algébriques, ni des propriétés dites ``\unsure{arithme|arithmé-}\newline
tiques'' de la cohomologie des variétés algébriques (ou encore, de la ``géométrie
des motifs''), aussi longtemps que la question de la validité de ces conjectures
n'est pas réglée. Aujourd'hui encore comme lors du Congrès de Bombay en 1968,
je considère cette question, avec celle de la résolution des singularités, \unsure{comme}
l'une des deux questions les plus fondamentales qui se posent en géométrie algé-\newline
brique. Je sens bien la portée de l'une et de l'autre! Cette fécondité \unsure{potentie|potentiel-}\newline
le ne pourra manquer de se manifester, dès l'instant où on ne se bornera plus à
contourner cahin-caha une conjecture décrétée ``trop difficile'', et où quelqu'un
prendra la peine enfin de retrousser ses manches et de s'y coltiner!
